% ================================================================================================ capítulo 3
\chapter{Medir con tu propio equipo}\label{cap:propio}

\begin{resumen}
\textbf{En este capítulo:} qué datos de tu telescopio, tu cámara y tu máscara necesita TurbuLENS, cómo agregarlos en
\boton{Dispositivos} para que queden guardados, cómo convertirlos en un perfil y cómo medir la escala de placa con el
cielo cuando no tienes los datos del fabricante. No hace falta ser del OAN ni tener sus equipos.
\end{resumen}

TurbuLENS trae los telescopios, las cámaras y las máscaras del Observatorio Astronómico Nacional porque son los que
usamos. Pero el método no depende de ellos: cualquier telescopio con una máscara de dos orificios y cualquier cámara que
grabe video sirven. Lo único que cambia es de dónde salen tres números: $D$, $B$ y la escala de placa.

\section{Lo que necesitas saber de tu equipo}

\begin{table}[H]
\centering\small
\begin{tabularx}{\linewidth}{@{}l Y Y@{}}
\toprule
Equipo & Dato & Dónde se encuentra\\
\midrule
Telescopio & Apertura y \textbf{distancia focal}, en mm. & En la etiqueta del tubo («D = 80 mm, F = 480 mm») o en el
manual. Con un reductor o una lente de Barlow, multiplica la focal por su factor (×0,63, ×2\ldots).\\
Cámara & \textbf{Ancho del sensor} en mm y en píxeles (o el tamaño del píxel en µm). & En la ficha técnica
(«sensor de 4,8 × 3,6 mm, 1280 × 960 px, píxel de 3,75 µm»).\\
 & Si es una cámara de fotos o un celular: el \textbf{ancho del video} en píxeles. & En la información del video
(\menu{Video}): 1920 en Full HD, 1280 en HD, 3840 en 4K.\\
Máscara & \textbf{$D$}, diámetro de cada orificio, y \textbf{$B$}, la distancia entre sus centros. & Mídelos con una regla
o un calibrador.\\
\bottomrule
\end{tabularx}
\caption{Los datos que hacen falta.}
\end{table}

\subsection{Hacer una máscara}

La máscara es un disco opaco (cartón paja, plástico negro, una tapa impresa en 3D) del diámetro de la boca del tubo, con
dos orificios circulares iguales. Unas reglas sencillas:
\begin{itemize}
\item Los dos orificios deben quedar \textbf{enteros dentro de la apertura} y, en un telescopio con espejo secundario
(Schmidt-Cassegrain, Newton), \textbf{fuera de su sombra}. Ponlos a lado y lado del centro, a la misma distancia.
\item $B$ debe ser al menos el doble de $D$; entre 3 y 5 veces es lo habitual. El modelo de Tokovinin que usa TurbuLENS
está pensado para orificios bien separados, y con ellos los dos puntos no se tocan al desenfocar.
\item $D$ entre 2,5 y 6~cm funciona bien con estrellas brillantes. Orificios más grandes dan más luz; más pequeños, puntos
más grandes y débiles.
\item Bordes limpios y sin rebabas: un borde irregular deforma los puntos.
\end{itemize}
Con la máscara puesta y el telescopio enfocado, las dos imágenes se superponen. Desenfoca un poco, hasta que se vean
\textbf{dos puntos separados} unas dos o tres veces su tamaño: esa es la imagen que mide TurbuLENS.

\section{Agregar tus equipos}\label{sec:agregar}

En \menu{Parámetros}, pulsa \botonp{Dispositivos (OAN y propios)…}. La ventana tiene una fila para el telescopio, otra
para la cámara y otra para la máscara, con los equipos del OAN y los tuyos.

\begin{capturam}[0.94\linewidth]{oan_propio}{La ventana \menu{Dispositivos} con un telescopio, una cámara y una máscara
propios elegidos (marcados con ★).}
\item Los \textbf{equipos}: los del OAN y los tuyos, marcados con ★. Haz clic en uno para elegirlo.
\item \boton{+ Propio…} agrega uno nuevo en esa fila.
\item \boton{Quitar} borra el equipo propio elegido (los del OAN no se pueden quitar).
\item \campo{Grabé en video} y el ancho del video, tomado del video cargado.
\item El \textbf{resultado}: la escala de placa, $D$, $B$, la focal, el píxel (efectivo, si es video) y una nota de cada equipo.
\item \botonp{Aplicar} pasa los valores a \menu{Parámetros}.
\end{capturam}

Con \boton{+ Propio…} se abre un formulario según la fila:

\begin{description}[style=nextline,leftmargin=1.2em,font=\bfseries\color{acentoosc}]
\item[Telescopio] Nombre, apertura (10 a 2000~mm) y distancia focal (50 a 20\,000~mm).
\item[Cámara] Nombre, ancho del sensor (1 a 60~mm) y ancho del sensor en píxeles. Si es una cámara de fotos o un
celular, marca \campo{El video sale con menos píxeles que el sensor} y escribe el ancho del video.
\item[Máscara] Nombre, $D$ y $B$. TurbuLENS no deja guardar una máscara con $B \le D$, porque los orificios se tocarían.
\end{description}

\begin{capturam}[0.62\linewidth]{dispositivo_propio}{El formulario de una cámara propia.}
\item El nombre con que aparecerá en la lista.
\item El ancho del sensor, en mm.
\item El ancho del sensor, en píxeles.
\item Márcala si es una cámara de fotos o un celular.
\item El ancho del video en píxeles.
\end{capturam}

Los equipos propios quedan guardados en la configuración de TurbuLENS de tu computador: aparecen la próxima vez que
abras el programa, en el modo medición y en el educativo.

\begin{nota}[¿Por qué solo el ancho?]
La escala de placa se calcula con el tamaño de un píxel, y para eso basta el ancho del sensor dividido entre su número
de píxeles. En las cámaras de fotos, ese mismo ancho se reparte entre los píxeles del video, que son menos: por eso
importa el ancho del video.
\end{nota}

\section{Un ejemplo completo}

Un club de astronomía mide con un refractor de 80~mm f/6 (focal de 480~mm), una cámara astronómica ZWO ASI120MM (sensor de
4,8~mm de ancho y 1280~px) y una máscara con orificios de 2~cm separados 7~cm:

\begin{enumerate}
\item En \boton{Dispositivos}, \boton{+ Propio…} en la fila del telescopio: «Refractor 80 mm f/6», apertura 80~mm, focal
480~mm.
\item \boton{+ Propio…} en la fila de la cámara: «ZWO ASI120MM», 4,8~mm y 1280~px. La casilla del video queda desmarcada:
una cámara astronómica graba con todos sus píxeles.
\item \boton{+ Propio…} en la fila de la máscara: «Máscara del club», $D = 2$~cm, $B = 7$~cm.
\item El resultado se ve al instante:
\[
p = \frac{4{,}8\ \text{mm}}{1280\ \text{px}} = 3{,}75\ \text{µm},\qquad
s = \frac{206{,}265 \times 3{,}75}{480} = 1{,}611\,''/\text{px}.
\]
\item \botonp{Aplicar}. En \menu{Parámetros}, \boton{Guardar…} en la fila del perfil: «Club: refractor 80 + ASI120».
\end{enumerate}

La próxima noche basta con elegir el perfil.

\begin{cuidado}[Celular en el ocular.]
Si grabas con un celular apoyado en el ocular (fotografía afocal), la escala depende del ocular, del celular y de su
zoom, y no sale de la focal del telescopio. En ese caso mídela con el cielo, como se explica a continuación, y escríbela a
mano en \menu{Parámetros}.
\end{cuidado}

\section{Medir la escala con el cielo}\label{sec:medir_escala}

Los datos de fábrica tienen tolerancias: una focal «de 2032~mm» puede ser de 2000 o de 2080 según dónde quede el foco,
y un reductor rara vez da exactamente su factor nominal. Si quieres el valor real de tu montaje, mídelo con una
estrella. Hay dos maneras.

\subsection{Con una estrella doble}

Graba, con el mismo montaje y la misma cámara, una estrella doble de separación conocida, \textbf{sin la máscara} y
enfocada. Albireo ($\beta$~Cyg, unos 34,4″) y Mizar A–B ($\zeta$~UMa, unos 14,4″) son fáciles; consulta el valor actual en
un catálogo de dobles, porque las separaciones cambian lentamente con la órbita. Luego:
\begin{enumerate}
\item Guarda un cuadro como imagen (en \menu{Parámetros}, \campo{Guardar los cuadros como imágenes}, en FITS o PNG) o
toma una foto con la misma cámara.
\item Ábrelo en el \menu{Visor de píxeles} del modo Imágenes, dibuja una región alrededor de la primera estrella y pulsa
\boton{Medir la estrella de la región}: anota su centro $(x_1, y_1)$. Repite con la segunda $(x_2, y_2)$.
\item La escala es
\[
s = \frac{\text{separación en el cielo}\ ('')}{\sqrt{(x_2-x_1)^2 + (y_2-y_1)^2}\ (\text{px})} .
\]
\end{enumerate}

\subsection{Con el movimiento del cielo}

Apaga el motor de seguimiento y deja que una estrella cruce la imagen: el cielo gira $15{,}041''$ por segundo de tiempo
multiplicado por $\cos\delta$, con $\delta$ la declinación de la estrella. Graba unos 30 segundos, procesa el video en el
modo medición con \campo{Seguir a la estrella} y exporta las mediciones en CSV. Las columnas \codigo{tiempo\_s}, \codigo{x1\_px} y \codigo{y1\_px} (la posición en el cuadro completo, no
en la región) dan la velocidad $v$ de la estrella en píxeles por
segundo; un ajuste lineal en una hoja de cálculo basta. Entonces
\[
s = \frac{15{,}041\,\cos\delta}{v}\quad[''/\text{px}] .
\]
Con una estrella cerca del ecuador celeste (Altair, Rigel, Procyon) el movimiento es más rápido y la medida más precisa.
Aquí puedes dejar la máscara puesta: los dos puntos se mueven igual.

\section{Antes de salir a medir}

\begin{pasos}[Lista de chequeo]
\begin{itemize}[label=$\square$,leftmargin=1.4em]
\item El perfil de tu montaje está guardado, o tienes a mano la focal, el sensor y $D$ y $B$.
\item La cámara graba al máximo de cuadros por segundo posible, con la exposición más corta que deje ver bien los dos
puntos y la ganancia (o el ISO) necesaria para que no se saturen.
\item Si es una cámara de fotos: sabes a qué resolución graba (1920 o 1280 de ancho).
\item Elegiste una estrella brillante y alta (más de 45° sobre el horizonte si puedes) y anotas la hora, para calcular su
altura después.
\item El telescopio está estable y lejos de paredes y techos calientes (el calor que sueltan es turbulencia local).
\item Grabas varios videos de 30 a 60 segundos en lugar de uno muy largo.
\end{itemize}
\end{pasos}
